% ─────────────────────────────────────────────────────────────────────────────
\section{Uniform Multi-Attack ERM Baseline}
\label{sec:ch4_multi_attack_erm}

Uniform Multi-Attack empirical risk minimization (ERM) serves as the main baseline
for this work. It trains on multiple source attacks within the same minibatch but
assigns equal weight to all generated adversarial examples. This makes it a strong
and simple reference point for evaluating whether adaptive reweighting provides
additional benefit.

Given a minibatch
\[
    \mathcal{B} = \{(\mathbf{x}_i,y_i)\}_{i=1}^{N},
\]
we partition it into \(M\) disjoint subsets
\[
    \mathcal{B}_1,\ldots,\mathcal{B}_M,
    \qquad
    \mathcal{B}_m \cap \mathcal{B}_{m'} = \emptyset
    \;\; \text{for } m\ne m',
    \qquad
    \bigcup_{m=1}^{M}\mathcal{B}_m = \mathcal{B}.
\]
Each subset \(\mathcal{B}_m\) is assigned to one source attack \(A_m\). For each
\((\mathbf{x}_i,y_i)\in\mathcal{B}_m\), we generate
\[
    \mathbf{x}_i^{\mathrm{adv}}
    =
    A_m(f_\theta,\mathbf{x}_i,y_i).
\]

The adversarial examples from all source attacks are then concatenated into one mixed
adversarial batch. The uniform multi-attack loss is
\begin{equation}
    \mathcal{L}_{\mathrm{ERM}}
    =
    \frac{1}{N}
    \sum_{i=1}^{N}
    \ell\!\left(
        f_\theta(\mathbf{x}_i^{\mathrm{adv}}),
        y_i
    \right).
    \label{eq:erm_loss}
\end{equation}

In our main setting, the source attacks are PGD-\(\ell_\infty\) and DDN-\(\ell_2\),
so the model is exposed to both \(\ell_\infty\) and \(\ell_2\) perturbation
geometries during training. However, the objective in
Equation~\eqref{eq:erm_loss} treats all adversarial examples equally. It does not
explicitly identify which attack, class, or latent subgroup is harder. This limitation
motivates the adaptive reweighting methods introduced in the following sections.
